\section{Generalized synthetic control: preserve the panel target across a wrapper}
\label{exec21:gsynth-scope}

A branch manager evaluates a staggered rollout of a servicing policy. Some
branches adopt early, some late, and some remain untreated. A useful analysis
must say whether its answer concerns a typical branch, a typical treated
branch-month, or a portfolio weighted by an explicitly chosen exposure.
These targets differ before any model is fitted. A package wrapper can change
the answer by losing a weight or substituting an estimator even when the
underlying numerical engine is otherwise appropriate.

The selected pre-cutoff \texttt{gsynth}~1.4.0 source is a compatibility
wrapper around \texttt{fect}. Its maintainers describe the wrapper as retained
solely for backward compatibility. The reviewed dependency candidate is
\texttt{fect}~2.4.5, declared on 29~May~2026 and published through the selected
CRAN revision on 30~May. Matching pinned upstream and CRAN wrapper, manual,
license and change-log files agree byte for byte. This establishes a concrete
source comparison; it does not establish an installed combination. The
wrapper's declared lower bound, \texttt{fect >= 2.0.0}, permits multiple engine
versions with potentially different behavior. \cite{exec21_gsynth_wrapper,exec21_gsynth_engine}

\subsection{Weights and exposure duration define the business quantity}
\label{exec21:gsynth-target}

Let $i$ index branches and $t$ calendar periods. Let $D_{it}$ indicate an
active intervention, $\mathcal S$ denote the retained treated cells, and
$\delta_{it}=Y_{it}(1)-Y_{it}(0)$ the corresponding effect in stated monetary
units. With positive prespecified weights $w_{it}$, two possible targets are
\begin{equation}
\label{exec21:gsynth-weights}
 \tau_{\rm cell}=\frac{1}{|\mathcal S|}\sum_{(i,t)\in\mathcal S}\delta_{it},
 \qquad
 \tau_w=\frac{\sum_{(i,t)\in\mathcal S}w_{it}\delta_{it}}
                  {\sum_{(i,t)\in\mathcal S}w_{it}}.
\end{equation}
If $\mathcal I$ is the retained treated-unit set and $T_i^+$ its observed
treated periods, a third target gives each unit equal weight:
\begin{equation}
\label{exec21:gsynth-units}
 \tau_{\rm unit}=\frac{1}{|\mathcal I|}\sum_{i\in\mathcal I}
       \frac{1}{T_i^+}\sum_{t:(i,t)\in\mathcal S}\delta_{it}.
\end{equation}
Identification of the missing $Y_{it}(0)$ is a separate question from these
definitions. A sample's retained cells need not represent the original
portfolio when observation or unit-removal rules depend on risk or outcome.

The selected \texttt{gsynth} manual promises that \texttt{weight} specifies
the weighting variable for the weighted average effect. The complete
95-line wrapper declares that argument but never reads it again and never
passes it to \texttt{fect}. Its documented weighted interface therefore has
a concrete source-level defect. A successful call or the stored call string
would not demonstrate that the requested weights reached estimation.
This is a source finding; the current review has not executed the native
wrapper. \cite{exec21_gsynth_wrapper}

For an exact illustration, two equally represented treated cells have effects
$-6$ and $10$ dollars. Their unweighted mean is $2$, whereas weights $9$ and
$1$ give $(-54+10)/10=-22/5=-4.4$. If the decision assumes a cost of USD~1
on the corresponding average-exposure basis, the two calculated margins are
USD~1 and USD~$-5.4$. A weighting omission can reverse the indicated decision;
these are authored arithmetic examples, not measured branch results or
observed outputs of either package. \cite{exec21_gsynth_reference}

Exposure duration creates another distinction. Suppose one unit experiences
effect $10$ in each of three treated periods, while another experiences
effect $-6$ in one treated period. Equation~\eqref{exec21:gsynth-weights}
gives $\tau_{\rm cell}=6$ and equation~\eqref{exec21:gsynth-units} gives
$\tau_{\rm unit}=2$. At an assumed per-period cost of $4$, averaging the
net effects using the same respective measures gives margins $2$ and $-2$.
A longer observed treatment history should not silently make one branch
three times as influential in an analysis intended to describe a typical
branch. \cite{exec21_gsynth_reference}

The selected engine exposes both treated-cell and unit-average calculations.
It also distinguishes weights for outcome-model fitting from weights for
aggregation. Supplying \texttt{W} fills both roles; supplying only
\texttt{W.agg} is intended to change aggregation, while \texttt{W.est}
alone leaves canonical aggregation unweighted. The inspected source rejects
distinct columns for the two roles, despite an error message referring to
earlier versions and a future release. Its executed conditional, rather than
that stale version message, determines the limitation. Zero-weight rows are
dropped, and weighted aggregates are routed into the ordinary output slots.
A proposed adapter must therefore record both the requested role and the
actual retained population; blindly multiplying supplied weights is not a
general substitute for missing role support. \cite{exec21_gsynth_engine}

\subsection{A factor model needs information in the relevant directions}
\label{exec21:gsynth-rank}

A useful working model for an untreated outcome is
\begin{equation}
\label{exec21:gsynth-factor}
 Y_{it}(0)=\alpha_i+\xi_t+x_{it}^{\mathsf T}\beta
                     +\lambda_i^{\mathsf T}f_t+\varepsilon_{it},
\end{equation}
where $x_{it}$ contains covariates with a justified intervention-time meaning,
$f_t$ represents common latent factors and $\lambda_i$ is the unit's loading.
Even if donor observations reveal $f_t$ exactly, a treated unit's pre-period
outcomes must determine the component of $\lambda_i$ that matters after
treatment. Covariates caused by the intervention, anticipation, spillovers
to donors, or changing loadings require additional arguments. A low prediction
error on past observations cannot establish those substantive conditions.

For a precise obstruction, suppress the intercepts, observed covariates and
errors. Let the factors in two pre-periods and one post-period be
\begin{equation}
\label{exec21:gsynth-witness}
 f_1=(1,0)^{\mathsf T},\quad f_2=(2,0)^{\mathsf T},\quad
 f_3=(1,1)^{\mathsf T}.
\end{equation}
Donor loadings $(0,1)$ and $(1,1)$ generate observed paths $(0,0,1)$ and
$(1,2,2)$. A treated loading of either $(2,0)$ or $(2,4)$ gives the same
pre-treatment path $(2,4)$. Hold its observed post-treatment outcome at $8$.
The two compatible untreated post outcomes are $2$ and $6$, giving effects
$6$ and $2$. With an assumed intervention cost of $4$, the net values are $2$
and $-2$. The complete observed panels agree between these two worlds.
\cite{exec21_gsynth_reference}

Writing $F_{\rm pre}$ for the pre-period factor matrix, the direction
$v=(0,1)^{\mathsf T}$ satisfies
\begin{equation}
\label{exec21:gsynth-null}
 F_{\rm pre}v=0,\qquad f_3^{\mathsf T}v=1.
\end{equation}
The pre-period factor rank is one while the full factor rank is two.
Repeating the two pre-factor rows over six distinct calendar periods does
not change that rank. A third donor with loading $(2,1)$ can raise the donor
count without resolving this missing treated-unit information. Thus minimum
period and donor counts are useful guards, but are not proofs of the relevant
rank condition. This example lies outside full pre-factor-rank qualification;
it is not evidence that an estimator valid under that qualification fails.
\cite{exec21_gsynth_reference}

The inspected engine includes count/rank restrictions and, in one unbounded
treated-loading branch, solves using the inverse of
$F_{\rm pre}^{\mathsf T}F_{\rm pre}$. A failed solve returns unavailable
effect fields. Near singularity can remain problematic even when the solve
returns. Its inspected rank-selection paths aggregate cross-validation
scores and can apply a minimum-error, one-standard-error or legacy
one-percent selection rule. The final rule may override an intermediate
rank update. These are predictive selection procedures, not tests proving
counterfactual identification. A restricted loading model, an informative
additional pre-period, an external loading restriction or a prospective
experiment could address a particular ambiguity, but each changes the
information or assumptions of the problem. \cite{exec21_gsynth_engine}

\subsection{Defaults, missingness and uncertainty must travel together}
\label{exec21:gsynth-inference}

In the selected wrapper, the default estimator is \texttt{gsynth}. Legacy
\texttt{EM=TRUE} selects \texttt{ife} only if the estimator argument was
omitted. The default requested inference is parametric for the gsynth route
and bootstrap for the IFE or matrix-completion routes. A parametric request
for the latter routes is changed to bootstrap with a warning; the legacy
name \texttt{nonparametric} maps to bootstrap. The wrapper also sets unit
fixed effects, five CV folds and tolerance $10^{-3}$, whereas the direct
engine's inspected defaults include two-way effects, twenty folds and
tolerance $10^{-5}$. Direct-engine adoption therefore requires explicit
configuration rather than assuming identical defaults.
\cite{exec21_gsynth_wrapper,exec21_gsynth_engine}

The engine's gsynth route forces a never-treated donor pool and rejects
treatment reversals. Too few pre-treatment observations or excessive
missingness can remove units; the returned result must disclose this change
in population. A package family's support for switching treatments does not
override the restrictions of this selected method. Missing-cell masks also
enter fitting and effect aggregation, so reporting only the number of
input rows conceals the denominator actually used. \cite{exec21_gsynth_engine}

For parametric uncertainty, the selected engine requires the isolated
never-treated pool. Its automatic error mode uses an autoregressive strategy
when cells are missing and an empirical strategy otherwise; explicitly
requested empirical or wild modes require a complete panel. In inspected
case-bootstrap code, units or cluster identifiers are resampled, with
support-conditioned retries in unit-resampling branches. Ordinary jackknife
with a cluster argument emits a warning that clustering is ignored. These
are different sampling mechanisms; a generic confidence-interval label
cannot make them interchangeable. \cite{exec21_gsynth_engine}

The inspected average-effect uncertainty uses the standard deviation of
bootstrap draws with missing values removed and, by default, a normal
interval. The basic-pivot option has a different construction. Any future
qualification must preserve attempted, valid and failed draws, resampling
units, selection/refitting behavior and reasons for unavailable output.
Dropping failed draws from a standard-deviation calculation is not evidence
that inference remains valid. Dependence across branches, common calendar
shocks, model misspecification and post-selection uncertainty each require a
declared law and a specific argument. This source comparison has executed
neither native package and makes no empirical coverage claim.
\cite{exec21_gsynth_engine}

\subsection{Reuse through an explicit panel contract}
\label{exec21:gsynth-integration}

The proposal retains the pinned wrapper as a compatibility reference and
favors evaluating a direct \texttt{fect} worker if a prototype is funded.
The ignored weight and implicit remapping make an unrestricted wrapper
pass-through unsuitable for a new interface. Porting the engine immediately
would duplicate panel conversion, rank selection, native numerical routines
and inference before a behavioral comparison exists. Classic synthetic
control, difference-in-differences, matrix-completion and state-space
approaches remain substantive comparators whose targets and identifying
assumptions must be matched to the banking question.

Both packages declare \texttt{MIT + file LICENSE}; the saved license files
contain the holder and year. That declaration supports identifying the
stated license, but this review has not inspected the complete standard
template or every dependency's terms. The wrapper says it needs no
compilation, while the selected engine requires compiled components,
R~4.1 or later, Rcpp and RcppArmadillo. A deployment estimate must include
that environment. Process isolation does not itself establish rights,
reproducibility or numerical validity. \cite{exec21_gsynth_wrapper,exec21_gsynth_engine}

A common library, CLI and MCP request should carry unique unit/time keys,
the complete calendar, action timing, outcome units, the donor population,
observation masks, target weights and their fitting/aggregation roles. It
should also freeze the method, fixed effects, rank candidates, selection
rule, tolerance, resampling law and dependency identities. The response
should return the actual configuration, retained sample and denominators,
conditioning and convergence diagnostics, original warnings, named effect
aggregates and a separately typed uncertainty result. An agent translating
the business question may propose these fields; it cannot infer a missing
weighting basis or certify donor validity from a fluent explanation.

Before admission, a public reference comparison should reproduce the
wrapper's weight omission and compare an explicitly configured engine on
known target calculations. It should then exercise unequal exposure
durations, missingness, reversals, donor shortages, weak factor support,
rank resets, failed solves, cluster warnings and bounded runtime failures.
Known-law factor panels can test selected calibration claims before any
costly bank experiment, provided selection and failure-denominator rules
are frozen in advance. Total cost includes validation, conversion, rank
search, repeated fits, uncertainty, serialization, review and maintenance.
These public comparisons can establish whether a panel prototype merits
investment while preserving the assumptions a later banking study must test.
\cite{exec21_gsynth_reference}
